\section{Discussion}
\subsection{시뮬레이션 장벽의 원인}
Section~V-D의 오류 분석은 시뮬레이션 장벽이 정보의 부재가 아니라 \emph{계산의 불안정성}에서 비롯됨을 시사한다. 설정 파일에는 OSPF cost, BGP neighbor, 인터페이스 IP와 같이 경로 계산에 필요한 정보가 이미 포함되어 있다. 그러나 이러한 정보를 정답으로 변환하려면 다수 노드에 걸친 SPF 계산과 VRF별 포워딩 상태 구성이 필요하며, 이는 단순한 텍스트 회상과는 다른 계산적 추론을 요구한다. 토큰 예측 방식의 LLM은 이 과정에서 중간 단계 오류를 누적시키며, 그 결과는 주로 빈 응답과 추론 회피, 또는 중간 홉을 잘못 삽입하는 추가 홉 오류의 형태로 나타난다. 이러한 실패 패턴이 파라미터 규모와 아키텍처에 관계없이 반복된다는 점은, 본 연구가 관찰한 한계가 개별 모델의 성능 차이라기보다 토큰 예측 패러다임 자체의 구조적 제약에 가깝다는 해석을 뒷받침한다. Mondal et al.~\cite{ref30}이 GPT-4에서도 라우터 설정 합성에 검증기 피드백 루프가 필요하다고 보고한 결과는 이러한 관찰과 맥을 같이한다. 따라서 네트워크 운영 관점에서 L1\textasciitilde L3 수준의 사실 추출과 정적 확인에는 LLM을 직접 활용할 수 있지만, L4/L5 수준의 동작 추론에는 정형 분석 도구의 병행이 사실상 불가피하다.

\subsection{구조 효과와 도구 효과의 분리}
기존 연구~\cite{ref14, ref15}는 Multi-Agent 도입에 따른 성능 향상을 보고하지만, 그 향상의 원천이 에이전트 구조인지 연결된 도구인지까지는 분리해 설명하지 않는다. 이 구분은 시스템 설계 관점에서 중요하다. 만약 성능 향상의 핵심이 구조적 분업이 아니라 도구 접근성에 있다면, 복잡한 다중 에이전트 아키텍처를 구축하는 것보다 더 단순한 구조에 적절한 네트워크 분석 도구를 연결하는 편이 실용적으로 더 타당할 수 있기 때문이다.

본 연구의 3-way 비교는 바로 이 지점을 겨냥한다. $\Delta$(구조)는 도구 없이 에이전트 구조만 추가했을 때의 효과를, $\Delta$(도구)는 동일한 구조에 도구를 통합했을 때의 추가 효과를 나타낸다. L1\textasciitilde L3에서는 질문 분해와 답변 통합 과정이 일부 이점을 제공할 수 있으나, L4/L5에서는 하위 질문 자체가 traceroute나 what-if 분석과 같은 시뮬레이션 절차를 필요로 한다. 따라서 구조만으로는 한계를 넘기 어렵고, 실제 성능 개선은 도구 호출과 그 결과 해석이 결합될 때 비로소 나타난다. 이 구간에서 LLM의 역할은 직접적으로 경로를 계산하는 주체라기보다, 어떤 도구를 호출할지 선택하고 그 결과를 사용자 질의에 맞게 해석하는 조정자로 재정의된다.

\subsection{파이프라인의 범용성과 실용적 가치}
이중 경로 파이프라인은 특정 토폴로지에 종속되지 않도록 설계되었다. 4개 Lab에 동일한 생성 절차를 적용한 결과, 코드 수정 없이 메트릭 커버리지가 52\%(Lab-A)에서 97\%(Lab-D)까지 확장되었고, QA/노드 비율도 84\textasciitilde126 범위에서 비교적 안정적으로 유지되었다. 이는 파이프라인이 단일 예제에 맞춘 일회성 구성이라기보다, 토폴로지 규모와 프로토콜 구성이 바뀌더라도 유사한 방식으로 적용될 수 있음을 보여준다. 실무적으로는 운영 조직이 자사 설정 파일만 제공하더라도 LLM 기반 도구의 사전 성능 평가, 변경 이후의 회귀 테스트, 또는 특정 운영 태스크에 대한 위험 점검에 이 프레임워크를 활용할 수 있다.

\subsection{타당성 위협}
\textbf{벤더 종속성.} 실험은 Cisco IOS로 한정하였다. 이는 벤더 간 구문 차이와 운영 관례의 영향을 배제하여 비교 조건을 단순화하기 위한 선택이다. Batfish가 벤더 독립적 데이터 모델을 사용하므로 파이프라인과 난이도 체계 자체는 Juniper JUNOS나 Arista EOS로 확장 가능하지만, 본 논문의 결과는 우선 Cisco IOS 환경에서 해석하는 것이 적절하다.

\textbf{Ground Truth 독립성.} L4/L5의 Ground Truth와 NetAlly의 핵심 추론 도구가 모두 Batfish에 의존하므로, NetAlly의 L4/L5 정확도는 부분적으로 ``올바른 Batfish API를 적절한 파라미터로 호출했는가''를 반영한다. 본 연구는 이러한 순환성을 통제하기 위해 세 가지 조치를 취하였다. 첫째, PNETLab 실장비 CLI와의 교차 검증으로 Batfish 출력의 신뢰성을 확인하였다(Section~III-E, Method~3). 둘째, 순환성에서 자유로운 L1\textasciitilde L3에서도 유사한 성능 격차 패턴이 나타나는지 함께 확인하였다. 셋째, 도구 호출 성공/실패별 TA-Acc를 분리 보고하여(Table~\ref{tab:tool-call-analysis}) ``API 호출 성공''과 ``추론 품질''을 구분하고자 하였다. 다만 Batfish 외의 독립 오라클, 예를 들어 실장비 기반 대규모 자동 검증을 통한 완전한 순환 해소는 향후 과제로 남아 있다.

\textbf{토폴로지 규모.} 본 연구의 최대 규모는 40노드로, 대규모 엔터프라이즈나 통신사업자 환경 전체를 대표하기에는 제한적이다. 그럼에도 QA/노드 비율이 비교적 안정적으로 유지되었고 Config Generator가 임의 규모의 토폴로지를 지원한다는 점은, 본 접근이 더 큰 환경으로 확장될 가능성을 보여준다. 다만 그 확장이 실제로 어떤 비용과 성능 저하를 수반하는지는 후속 실험으로 확인할 필요가 있다.